% !PRINTED 604
\weberpage{604}

%<<CONT
endlich. Die erste Summe in (9) ist

\weq{10}{\sum\frac{1}{\nu^s} = \frac{1}{s-1} + C_s,}

wo $C_s$ für $s = 1$ endlich bleibt, nämlich in den \gesperrt{Wert der Eulerschen Konstante} $0{,}5772\dots$ übergeht (Bd.\ II, S.\ 723).

Berücksichtigt man dies, so ergibt sich aus (9) das Resultat:

5. \gesperrt{Hat $S_i$ die Bedeutung (3), so ist}

\weq{11}{S_i = \frac{g}{s-1} + C_i,}

\gesperrt{worin $g$ eine von der besonderen Klasse $\overline{A}_i$ unabhängige von Null verschiedene Konstante ist, und $C_i$ eine von $\overline{A}_i$ abhängige Funktion von $s$, die für $s = 1$ endlich bleibt.}

Die Gruppe $h$ten Grades der Idealklasse $\overline{A}_i$ ist eine \gesperrt{Abelsche}, und es gibt daher $h$ Charaktere $\chi_\kappa(\overline{A}_i)$, die sämtlich $h$te Einheitswurzeln sind, darunter der \gesperrt{Hauptcharakter} $\chi_1$, der für jede Klasse $= 1$ ist. Ebenso hat für die Hauptklasse $\overline{A}_1$ jeder der Charaktere den Wert $+1$. Für diese Charaktere bestehen die Sätze:

\weq{12}{\begin{aligned}
&\sum^{i}\chi_\kappa(\overline{A}_i) = 0, \quad \sum^{\kappa}\chi_\kappa(\overline{A}_i) = 0, \\
&\kappa = 2,3,\dots,h, \quad i = 2,3,\dots,h, \\
&\sum^{i}\chi_1(\overline{A}_i) = h, \quad \sum^{\kappa}\chi_\kappa(\overline{A}_1) = h
\end{aligned}}

(Bd.\ II, \S\,13). Wenn \gesperrt{das ganze Ideal} $\mathfrak{a}$ in $\overline{O}$ in die Klasse $\overline{A}$ gehört, so setzen wir für jedes $\chi$:

\weq{13}{\chi(\mathfrak{a}) = \chi(\overline{A}),}

\gesperrt{und haben so für jedes Ideal in $\overline{O}$ die Charaktere nach der Gruppe $A$} bestimmt. Für irgend zwei Ideale $\mathfrak{a}$ und $\mathfrak{b}$ aus $\overline{O}$ besteht die Relation:

\weq{14}{\chi(\mathfrak{a})\,\chi(\mathfrak{b}) = \chi(\mathfrak{ab}).}

Nun bilden wir aus $S_i$ die $h$ Summe:
% [source sic: the printed singular "Summe" conflicts with the h sums Q_\kappa that follow]

\weq{15}{Q_\kappa = \sum^{i}\chi_\kappa(\overline{A}_i)\,S_i,}

wofür wir nach (3) und (13) auch setzen können:

\weq{16}{Q = \sum^{\mathfrak{a}}\frac{\chi(\mathfrak{a})}{(N\mathfrak{a})^s}.}
